%% ============================================================
%% Chapter 9: The Discrete Toy Model
%% ============================================================

\chapter{The Discrete Toy Model}
\label{ch:ch09-discrete-toy-model}
\section{An Explicit Toy Model of Holonomy}

Chapter~\ref{ch:ch06-distinction-holonomy} developed connection, curvature, and holonomy abstractly, and Theorem~\ref{thm:non-identifiability} asserted that identity becomes necessarily historical whenever a continuation connection possesses nonzero curvature or nontrivial holonomy, provided such connections actually exist. This short chapter discharges that existence claim with a fully worked, checkable example, small enough to verify by hand and general enough to make the abstract machinery of the preceding chapter concrete.

Consider the generators
\[
X = \begin{pmatrix} 0 & 1 \\ 0 & 0 \end{pmatrix}, \qquad Y = \begin{pmatrix} 0 & 0 \\ 1 & 0 \end{pmatrix}, \qquad [X,Y] = XY - YX = \begin{pmatrix} 1 & 0 \\ 0 & -1 \end{pmatrix} \neq 0.
\]
The group commutator $\gamma_\epsilon = e^{\epsilon X} e^{\epsilon Y} e^{-\epsilon X} e^{-\epsilon Y}$ satisfies, by the Baker--Campbell--Hausdorff formula,
\[
\gamma_\epsilon = I + \epsilon^2 [X,Y] + O(\epsilon^3).
\]
This algebraic fact alone is not yet a holonomy model, since no base manifold, bundle, connection, or loop in the base has been specified. To make it one, take $M = \mathbb{R}^2$, the trivial rank-two bundle, and the constant connection
\[
A = X\, dx + Y\, dy.
\]
Then $F = dA + A \wedge A = [X,Y]\, dx \wedge dy \neq 0$, and parallel transport around a small coordinate rectangle of side $\epsilon$ reproduces $\gamma_\epsilon$ above, up to a sign fixed by the transport convention and loop orientation. A system transported around this loop returns to its original base point $(0,0)$ but not to its original fiber state:
\[
\Delta_\gamma = \gamma_\epsilon - I = \epsilon^2 [X,Y] + O(\epsilon^3) \neq 0.
\]

No storage variable and no explicit historical record has been introduced anywhere in this construction. The memory is encoded entirely in the failure of transport to be path-independent. This is worth pausing on, because it is easy to read the abstract claims of Chapter~\ref{ch:ch06-distinction-holonomy} as a metaphor for memory rather than a literal mechanism for it: here, the base point $(0,0)$ genuinely returns to itself, the system genuinely traverses a closed loop, and yet the fiber state at the end of the traversal genuinely differs from the fiber state at the beginning, by an amount computable in closed form, $\epsilon^2[X,Y]$, with no register, log, or symbol anywhere storing that difference as data. The difference simply is the noncommutativity of the two generators, made geometric by the choice of connection.

The toy model also exhibits, explicitly, the pair of paths that Theorem~\ref{thm:non-identifiability} required to exist. The two orderings $e^{\epsilon X} e^{\epsilon Y}$ and $e^{\epsilon Y} e^{\epsilon X}$ around this loop realize exactly the pair $(\gamma_A, \gamma_B)$ the non-identifiability theorem needs: both orderings begin and end at the same base point, both are compatible with the same instantaneous observations at each intermediate point if only the base coordinate is observed, and yet the two orderings induce different fiber transformations whenever $[X,Y] \neq 0$. An observer restricted to the base manifold, watching only where the system is rather than what it carries in its fiber, cannot distinguish the two histories at any single instant. Only the accumulated holonomy, the full traversal, reveals that they differ.

This is the smallest possible demonstration of a claim that recurs, in less formal dress, throughout this monograph: that two processes can be observationally identical at every instant while remaining genuinely different in what they make admissible next. The toy model does not prove that any specific real system, a persona ensemble, an incident ledger, a paracosm, an institution, exhibits nonzero curvature of this particular kind; it proves only that the mathematical possibility Theorem~\ref{thm:non-identifiability} depends on is not vacuous, that connections with exactly this property exist, are computable by hand from a single matrix commutator, and are not an artifact of the abstract formalism's generality. Whatever curvature or holonomy a given real system possesses remains an empirical question about that system; what this chapter establishes is that the question is well-posed, because the mathematical structure it presupposes is not empty.
